\chapter{Propagation speed and kernel estimates}
\source{cheeger-gromov-taylor:page-0020,cheeger-gromov-taylor:page-0021,cheeger-gromov-taylor:page-0022,cheeger-gromov-taylor:page-0023,cheeger-gromov-taylor:page-0024,cheeger-gromov-taylor:page-0025,cheeger-gromov-taylor:page-0026,cheeger-gromov-taylor:page-0027}

The estimate (2.13) on $k_f(x_1,x_2)$ for $f\in \mathscr{F}_2^{\varphi,A}$ holds with $A=d/2$, and we have

\begin{equation}\tag{3.4}
|k_f(x_1,x_2)|\le c_n^2\left(\frac{2}{a}\right)^n \psi(r)^2,
\end{equation}

if

\begin{equation}\tag{3.5}
\rho(x_1,x_2)=a+r.
\end{equation}

If we do not assume a lower bound on the Ricci curvature, methods of estimating the gradient of $k_f(x_1,x_2)$ do not apply. We note that the assumption of $C^0$-bounded geometry implies a H\"older bound on $k_f(x_1,x_2)$ which improves the pointwise bound (3.4) for $f\in \mathscr{F}_2^{\varphi,A}$.

Indeed, according to Corollary~\ref{cor:derivative-tail-estimate} if $u$ is supported on $\Ba(x_2)$, then

\begin{equation}\tag{3.6}
\bigl\|\Delta^k f(\sqrt{-\Delta}) \Delta^l u\bigr\|_{M\setminus B_{r+a}(x_2)}
\le \frac{1}{\pi} C(C(2k+2l))^{2k+2l}\psi(r)\|u\|.
\end{equation}

Thus if $\rho(x_1,x_2)=r+2a$, then

\begin{equation}\tag{3.7}
\bigl\|\Delta^k f(\sqrt{-\Delta}) \Delta^l u\bigr\|_{B_a(x_1)}
\le \frac{C}{\pi}(C(2k+2l))^{2k+2l}\psi(r)\|u\|.
\end{equation}

We now argue similarly to the proof of Theorem~\ref{thm:kernel-estimates-ricci-lower-bound}. Given $u$ supported in $B_a(x_2)$, if $|y|\le B=\tfrac12 c^{-1}$, the power series

\begin{equation}\tag{3.8}
v(y,x)=\sum_{k=0}^{\infty}\frac{y^{2k}}{(2k)!}\,\Delta^k f(\sqrt{-\Delta})u
\end{equation}

converges to an element of $L^2([-B,B]\times B_a(x_1))$, and we have

\begin{equation}\tag{3.9}
\left(\frac{\partial^2}{\partial y^2}+\Delta\right)v=0,
\end{equation}

\begin{equation}\tag{3.10}
\|v\|_{[-B,B]\times B_a(x_1)}\le C_1\psi(r)\|u\|_{L^2}.
\end{equation}

Now in normal coordinates on $B_a(x_1)$ we have

\begin{equation}\tag{3.11}
\Delta w(x)=\sum_{j,k} g^{-1/2}\frac{\partial}{\partial x_j}
\left(g^{1/2}g^{ij}\frac{\partial w}{\partial x_k}\right),
\end{equation}

and our hypothesis that $M$ has $C^0$-bounded geometry precisely gives uniform ellipticity of (3.11), and hence of (3.9), with $C^0$ bounds on the coefficients. Thus the De Giorgi--Nash--Moser theorems on such elliptic equations in divergence form (see Gilbarg and Trudinger \cite{cgt-gilbarg-trudinger}) imply the pointwise estimate

\begin{equation}\tag{3.12}
|v(x_1)|\le C_2\psi(r)\|u\|
\end{equation}

for $u$ supported in $\Ba(x_2)$. Equivalently,
$f(\sqrt{-\Delta})\delta_{x_1}$ is $L^2$ on $B_a(x_2)$ and
\begin{equation}\tag{3.13}
 \bigl\|f(\sqrt{-\Delta})\delta_{x_1}\bigr\|_{B_a(x_2)}\le c_2\psi(r).
\end{equation}

Similarly,
\begin{equation}\tag{3.14}
 \bigl\|\Delta^l f(\sqrt{-\Delta})\delta_{x_1}\bigr\|_{B_a(x_2)}\le c_2(2Cl)^{2l}\psi(r).
\end{equation}

for $|y|\le B$, define
\begin{equation}\tag{3.15}
 v_{x_1}(y,x)=\sum_{l=0}^{\infty}\frac{y^{2l}}{(2l)!}\Delta^l f(\sqrt{-\Delta})\delta_{x_1}
\end{equation}
Then
\begin{equation}\tag{3.16}
 \left(\frac{\partial^2}{\partial y^2}+\Delta\right)v_{x_1}=0,
\end{equation}
\begin{equation}\tag{3.17}
 \bigl\|v_{x_1}\bigr\|_{[-B,B]\times B_a(x_2)}\le c_3\psi(r).
\end{equation}

The De Giorgi--Nash--Moser estimate applied again gives
\begin{equation}\tag{3.18}
 k_f(x_1,x_2)\le c_4\psi(r).
\end{equation}

and in fact yields the H\"older bound
\begin{equation}\tag{3.19}
 \bigl\|k_f(x_1,\cdot)\bigr\|_{C^\alpha(B_{a/2}(x_2))}\le C_5\psi(r).
\end{equation}

for some $\alpha>0$ depending only on the $C^0$ bound and the ellipticity constant in (3.11).

\begin{theorem}[Kernel estimates for $C^0$-bounded geometry]
\label{thm:c0-bounded-geometry-kernel-estimates}
\source{cheeger-gromov-taylor:page-0020,cheeger-gromov-taylor:page-0021}
\uses{cor:derivative-tail-estimate}
\dcref{Theorem 3.1}
\group{section-3}


If $M$ has $C^0$-bounded geometry and $f\in \mathcal{F}_2^{\varphi,A}$, then the kernel $k_f(x_1,x_2)$ satisfies (3.19) and
\begin{equation}\tag{3.20}
 \bigl\|k_f(\cdot,x_2)\bigr\|_{C^\alpha(B_{a/2}(x_1))}\le C_5\psi(r).
\end{equation}
\end{theorem}
\begin{proof}
\sketch Apply the De Giorgi--Nash--Moser estimate to (3.15)--(3.17), obtaining (3.18) and (3.19); symmetry of the kernel gives (3.20).
\end{proof}

The symmetry of $k_f(x_1,x_2)$ gives (3.20). If $f(\lambda)$ is the restriction to
$\lambda\in\mathbf R$ of a function $f(z)$ holomorphic on
\begin{equation}\tag{3.21}
 \Omega=\{z\in \mathbf{C}:\ |\operatorname{Im}z|\le W+\epsilon\}\cup\{z\in \mathbf{C}:\ |\operatorname{Im}z|\le B\,|\operatorname{Re}z|\},
\end{equation}
then $f\in \mathcal F_2^{\varphi,A}$ with $\varphi(s)=Ce^{-Ws}$, provided
\begin{equation}\tag{3.22}
 |f(z)|\le C(1+|z|)^m,\qquad z\in\Omega.
\end{equation}

This gives the following decay estimate.

\begin{corollary}[Exponential kernel decay under holomorphy]
\label{cor:c0-bounded-geometry-exponential-decay}
\source{cheeger-gromov-taylor:page-0021,cheeger-gromov-taylor:page-0022}
\uses{thm:c0-bounded-geometry-kernel-estimates}
\dcref{Corollary 3.2}
\group{section-3}


If $M$ has $C^0$-bounded geometry and $f(z)$ is holomorphic on (3.21) satisfying (3.22), then
\begin{equation}\tag{3.23}
 |k_f(x_1,x_2)|\le Ce^{-Wr},
\end{equation}
\end{corollary}
\begin{proof}
\sketch The holomorphy hypothesis places $f$ in $\mathcal F_2^{\varphi,A}$ with $\varphi(r)=Ce^{-Wr}$; apply the preceding theorem.
\end{proof}

provided
\begin{equation}\tag{3.24}
\rho(x_1,x_2)=r+2a.
\end{equation}

Assume now that $M$ has $C^\infty$-bounded geometry. The bounded-curvature estimates
control $k_f(x_1,x_2)$ for $\rho(x_1,x_2)\ge a$; near the diagonal we use the
pseudodifferential representation below. Let $f(\lambda)\in S_{\rho,0}^m(\mathbf R)$, i.e.,

\begin{equation}\tag{3.25}
|f^{(j)}(\lambda)|\le C_j(1+|\lambda|)^{m-j\rho}.
\end{equation}

We use the standard symbol classes and pseudodifferential calculus; see
\cite[Chapter~II]{cgt-taylor-pdo}. From the representation

\begin{equation}\tag{3.26}
f(\sqrt{-\Delta})u=\frac{1}{2\pi}\int_{-\infty}^{\infty}\hat f(s)\cos s\sqrt{-\Delta}\,u\,ds,
\end{equation}

use a partition of unity $1=\phi_1(s)+\phi_2(s)$, $\phi_1(s)\in C_0^\infty(-a,a)$, even, $\phi_1(s)=1$ for $|s|\le a/2$, to get

\begin{align}
f(\sqrt{-\Delta})u
&=\frac{1}{2\pi}\int_{-a}^{a}\phi_1(s)\hat f(s)\cos s\sqrt{-\Delta}\,u\,ds \notag\\
&\quad +\frac{1}{2\pi}\int_{-\infty}^{\infty}\phi_2(s)\hat f(s)\cos s\sqrt{-\Delta}\,u\,ds \notag\\
&=F_1u+F_2u.
\tag{3.27}
\end{align}

Note that

\begin{equation}\tag{3.28}
F_2u=(1-\Delta)^{-k}\int_{-\infty}^{\infty}
\left(1-\frac{d^2}{ds^2}\right)^k(\phi_2(s)\hat f(s))\cos s\sqrt{-\Delta}\,u\,ds
=(1-\Delta)^{-k}G_ku,
\end{equation}

with $\|G_ku\|\le C_k\|u\|$. Thus the standard elliptic estimates and the Sobolev embedding theorem give

\begin{equation}\tag{3.29}
\|F_2u\|_{C^l(B_a(x_1))}\le C_l\|u\|_{L^2}.
\end{equation}

Furthermore, if $u$ is supported in $B_a(x_1)$,

\begin{equation}\tag{3.30}
\|F_2u\|_{C^l(B_a(x_1))}\le C_{lk}\|u\|_{H^{-k}(B_a(x_1))},
\end{equation}

where the norm on the right is a Sobolev norm.

For the near-time term,
\begin{equation}\tag{3.31}
F_1u=\frac{1}{2\pi}\int_{-a}^{a}\varphi_1(s)\hat f(s)\cos s\sqrt{-\Delta}\,u\,ds,
\end{equation}
with $u$ supported in $B_a(x_1)$. For $|s|\le a$, replace
$\cos s\sqrt{-\Delta}$ by the standard geometrical-optics parametrix.

By shrinking $a$ if necessary, we can assume the parametrix for $\cos s\sqrt{-\Delta}$ is of the following form, in a normal coordinate system centered at $x_1$:
\begin{equation}\tag{3.32}
\cos s\sqrt{-\Delta}\,u(x)=\sum_{j=1}^{2}\int a_j(s,x,\xi)e^{i\varphi_j(s,x,\xi)}\hat u(\xi)\,d\xi.
\end{equation}
Here $\varphi_j$ solves the eikonal equation
\begin{equation}\tag{3.33}
\frac{\partial}{\partial s}\varphi_j = (-1)^j\lambda_1^j(x,d\varphi_j)=(-1)^j(d\varphi_j,d\varphi_j)^{\frac12},
\end{equation}
\begin{equation}\tag{3.34}
\varphi_j(0,x,\xi)=x\cdot \xi
\end{equation}
for $|s|\le a$, $x\in B_{2a}(x_1)$; $a_j(s,x,\xi)$ is determined by the usual transport equations of geometrical optics (see for example, \cite[Chapter~VIII]{cgt-taylor-pdo}), with $a_j(0,x,\xi)=\frac12$. Now if $u$ is supported in $B_a(x_1)$, then, for $|s|\le a$, $\cos s\sqrt{-\Delta}\,u$ is supported in $B_{2a}(x_1)$. If we set $\hat g(s)=\varphi_2(s)\hat f(s)$, then $g(\lambda)\in S^{m}_{\rho,0}$ differs from $f(\lambda)$ by an element of the Schwartz space $\mathcal S(\mathbb R)$ of rapidly decreasing functions. Thus we have
\begin{equation*}
F_1u=\sum_{j=1}^{2}\int_{-a}^{a}\hat g(s)a_j(s,x,\xi)e^{i\varphi_j(s,x,\xi)}\hat u(\xi)\,d\xi\,ds+S_1u
\end{equation*}
\begin{equation}\tag{3.35}
=\sum_j\int g(D_s)\!\left(a_j(s,x,\xi)e^{i\varphi_j(s,x,\xi)}\right)\Big|_{s=0}\hat u(\xi)\,d\xi+S_1u
\end{equation}
\begin{equation*}
=\int b(x,\xi)e^{ix\cdot \xi}\hat u(\xi)+S_1u,
\end{equation*}
where $S_1:\mathcal E'(B_a(x_1))\to C^\infty(B_{3a}(x_1))$, and $b(x,\xi)$ is given by the fundamental asymptotic expansion lemma for pseudodifferential operators. Thus
\begin{equation}\tag{3.36}
b(x,\xi)=\sum_{j=1}^{2}b_j(0,x,\xi),
\end{equation}
where
\begin{equation}\tag{3.37}
b_j(s,x,\xi)=e^{-i\varphi_j(s,x,\xi)}g(D_s)\!\left(a_j(s,x,\xi)e^{i\varphi_j(s,x,\xi)}\right).
\end{equation}

belongs to $S^m_{\rho,1-\rho}$ if $\rho \ge \frac12$ with
\begin{equation}\tag{3.38}
b_j(s,x,\xi)\sim a_j(s,x,\xi)g\!\left(\frac{\partial}{\partial s}\varphi_j\right)+\cdots.
\end{equation}
Hence using (3.33) and (3.34) we obtain
\begin{equation}\tag{3.39}
b(x,\xi)\sim g(\lambda_1(x,\xi))+\cdots.
\end{equation}
The remainder terms in (3.38) and (3.39) sum to elements of $S^{m-(2\rho-1)}_{\rho,1-\rho}$ for $\rho \ge \frac12$ and we have a complete asymptotic expansion for $\rho>\frac12$.

Denoting by $\OPS^m_{\rho,\delta}$ the set of pseudodifferential operators with symbols in the class $S^m_{\rho,\delta}$, we have the following conclusion.

\begin{theorem}[Pseudodifferential representation of $f(\sqrt{-\Delta})$]
\label{thm:pdo-representation-of-functional-calculus}
\source{cheeger-gromov-taylor:page-0023,cheeger-gromov-taylor:page-0024}
\dcref{Theorem 3.3}
\group{section-3}

Assume $M$ has $C^\infty$-bounded geometry. Let $f(\lambda)\in S^m_\rho(\R)$ be even with $\rho\ge\frac12$. If $u$ is supported in $\Ba(p)$ in normal coordinates centered at $p$, then
\begin{equation}\tag{3.40}
f(\sqrt{-\Delta})u=B_p(x,D)u;\ x\in \Ba(0),
\end{equation}
where $B_p(x,D)\in \OPS^m_{\rho,1-\rho}$ and
\begin{equation}\tag{3.41}
\{B_p(x,D): p\in M\}\text{ is bounded in }\OPS^m_{\rho,1-\rho}.
\end{equation}
If $\rho>\frac12$, its symbol has the complete asymptotic expansion (3.39), uniformly in $p$.
\end{theorem}
\begin{proof}
\sketch Split (3.26) at scale $a$, use finite propagation for the far-time term, and apply the geometric-optics parametrix (3.32)--(3.39) to the near-time term.
\end{proof}

Theorem~\ref{thm:pdo-representation-of-functional-calculus}, combined with the
off-diagonal kernel estimates above, gives $L^p$ bounds for the following
holomorphic symbol class.

\begin{definition}[Holomorphic symbol class on a strip]
\label{def:holomorphic-strip-symbol-class}
\source{cheeger-gromov-taylor:page-0024}
\dcref{Definition 3.1}
\group{section-3}

We say $f\in S_w^m$ if $f(z)$ is holomorphic on the strip
\begin{equation}\tag{3.42}
\Omega_w=\{z\in \C : |\Impart z|<W\},
\end{equation}
and satisfies the $S^m_{1,0}$ symbol estimates on $\Omega_w$:
\begin{equation}\tag{3.43}
|f^{(j)}(z)|\le C_j(1+|z|)^{m-j},\qquad z\in \Omega_w.
\end{equation}
\end{definition}

We have, according to (3.27),
\begin{equation}\tag{3.44}
f(\sqrt{-\Delta})=\Phi_1+\Phi_2,
\end{equation}
where the kernel of $\Phi_1$ is supported in the neighborhood $\{(x_1,x_2): \rho(x_1,x_2)\le 2a\}$ of the diagonal in $M\times M$ and is given in local normal coordinates as a bounded family of pseudodifferential operators. The kernel $\Phi_2(x_1,x_2)$ is smooth on $M\times M$ and satisfies the estimate
\begin{equation}\tag{3.45}
|\Phi_2(x_1,x_2)|\le Ce^{-Wr},\qquad r=\rho(x_1,x_2).
\end{equation}

From this we can deduce $L^{p}$ continuity results for $f(\sqrt{-\Delta})$, as follows. First we show $\Phi_1: L^{p}(M)\to L^{p}(M)$, $1<p<\infty$. Let $u\in L^{p}(M)$. We can write

\begin{equation}\tag{3.46}
u=\sum_{j=0}^{\infty}u_{j}
\end{equation}

where each $u_{j}$ is supported in a ball $B_{a}(p_{j})$ of radius $a$ in $M$, and, because we are assuming $M$ has bounded geometry, we can suppose there is a constant $K$ such that, for any $j$, $B_{2a}(p_{j})$ intersects no more than $K$ other balls $B_{2a}(p_{k})$, $p_{k}\neq p_{j}$. Thus we can arrange that

\begin{equation}\tag{3.47}
C^{-1}\sum_{j}\|u_{j}\|_{L^{p}}\leq \|u\|_{L^{p}}\leq C\sum_{j}\|u_{j}\|_{L^{p}}.
\end{equation}

Now $\Phi_{1}u=\sum_{j}\Phi_{1}u_{j}$, and operators $b(x,D)$ with symbols in $S^{0}_{1,0}$ are continuous on $L^{p}$, $1<p<\infty$, and bounded families of such operators have uniformly bounded operator norm on $L^{p}$, $1<p<\infty$; see \cite[Chapter~XII]{cgt-taylor-pdo}. Thus from (3.40) and (3.41) it follows that

\begin{equation}\tag{3.48}
\|\Phi_{1}u_{j}\|_{L^{p}}\leq C_{1}\|u_{j}\|_{L^{p}},
\end{equation}

provided $f\in S^{0}_{1,0}(\mathbf{R})$. Furthermore, $\Phi_{1}u_{j}$ is supported in $B_{2a}(p_{j})$. It follows that

\begin{equation}\tag{3.49}
\|\Phi_{1}u_{j}\|_{L^{p}}\leq C_{2}(p)\|u\|_{L^{p}}, \qquad \text{if } f\in S^{0}_{1,0}(\mathbf{R}).
\end{equation}

As for the $L^{p}$ continuity of $\Phi_{2}$, we use the estimate (3.45) together with the elementary fact that the $L^{p}$ operator norm of $\Phi_{2}$ is bounded by

\begin{equation}\tag{3.50}
\sup_{q}\int_{M}|\Phi_{2}(p,q)|\,d\vol(q)+\sup_{q}\int_{M}|\Phi_{2}(p,q)|\,d\vol\,p.
\end{equation}

The $L^p$ bound for the remainder follows from the pseudodifferential
calculus \cite[Chapter~XIII]{cgt-taylor-pdo}. If $M$ has
$C^{\infty}$-bounded geometry, or more generally Ricci curvature bounded below,
we have

\begin{equation}\tag{3.51}
\vol\{x\in M:\rho(x,p)\leq r\}\leq ce^{Kr}
\end{equation}

for some $K$ independent of $p$. Consequently, if (3.45) holds with $W>K$, we have a uniform bound on (3.50) and hence an $L^{p}$-operator bound on $\Phi_{2}$, $1\leq p\leq \infty$. Thus we have proved the following result.

\begin{theorem}[\texorpdfstring{$L^p$}{Lp} continuity for strip-holomorphic symbols]
\label{thm:lp-continuity-strip-holomorphic-symbols}
\source{cheeger-gromov-taylor:page-0025}
\dcref{Theorem 3.4}
\group{section-3}

Assume $M$ has $C^{\infty}$-bounded geometry. If
\begin{equation}\tag{3.52}
f\in \Swo, \qquad W>K,
\end{equation}
where (3.51) holds for $K$, then
\begin{equation}\tag{3.53}
f(\sqrt{-\Delta}): L^{p}(M)\to L^{p}(M), \qquad 1<p<\infty.
\end{equation}
\end{theorem}
\begin{proof}
\sketch The local part is uniformly $L^p$-bounded by (3.40)--(3.41); the decay (3.45) and volume growth (3.51), with $W>K$, give Schur bounds for the remainder.
\end{proof}
